\section{The Framework}\label{sec:framework}

This section lays out the FLM framework as a practitioner would encounter it: we describe the structural model, define target functionals, explain why naive neural network inference fails, present the influence function correction, and classify problems into three ``regimes'' that determine how a key quantity, the expected Hessian $\Lambda(x)$, is handled.

\subsection{Setup}\label{sec:setup}

We observe $n$ i.i.d.\ observations $(Y_i, T_i, X_i)$ where:
\begin{itemize}
    \item $Y_i$ is the outcome (chosen product, sales quantity, duration, etc.)
    \item $T_i$ is the treatment or ``action'' variable (price, product attributes, dosage)
    \item $X_i \in \mathbb{R}^{d_x}$ are individual covariates that drive heterogeneity
\end{itemize}

A \emph{structural model} specifies a loss function $\ell(y, t, \theta)$ parameterized by structural parameters $\theta \in \mathbb{R}^{d_\theta}$. The true parameters for individual $i$ solve:
\begin{equation}\label{eq:theta_star}
    \theta^*(X_i) = \arg\min_\theta \; E[\ell(Y, T, \theta) \mid X = X_i]
\end{equation}

\paragraph{Concrete example.} In our H\&M application:
\begin{itemize}
    \item $Y_i \in \{0, 1, \ldots, J-1\}$ is the chosen product from a set of $J$ alternatives
    \item $T_i \in \mathbb{R}^{J \times K}$ contains the attributes of each alternative (log-price, style embeddings)
    \item $X_i \in \mathbb{R}^{64}$ is the consumer's learned embedding
    \item $\theta^*(X_i) = (\beta_{\text{price}}(X_i), \beta_{\text{style},1}(X_i), \ldots, \beta_{\text{style},5}(X_i))$ are the consumer's taste parameters
    \item The loss is the multinomial logit negative log-likelihood:
\begin{equation}
    \ell(y, t, \theta) = -V_{y} + \log \sum_{j=0}^{J-1} e^{V_j}, \quad V_j = x_j' \theta
\end{equation}
\end{itemize}

The framework is \emph{not} limited to discrete choice. Table~\ref{tab:models} lists the structural models currently available in the \texttt{deep-inference} package, along with their loss functions and parameter dimensions.

\begin{table}[ht]
\centering
\caption{Available structural models in the \texttt{deep-inference} package.}\label{tab:models}
\begin{tabular}{llll}
\toprule
\textbf{Model} & \textbf{Loss} $\ell(y,t,\theta)$ & \textbf{Link} & $d_\theta$ \\
\midrule
Linear      & $(y - \alpha - \beta t)^2$                                & Identity & 2 \\
Logit       & $\log(1+e^\eta) - y\eta$                                  & Logistic & 2 \\
Poisson     & $e^\eta - y\eta$                                          & Log      & 2 \\
Gamma       & $y/\mu + \log \mu$                                        & Log      & 2 \\
Multinomial & $-V_y + \log\sum_j e^{V_j}$                               & Softmax  & $(J\!-\!1)\!+\!K$ \\
Custom      & Any differentiable $\ell(y,t,\theta)$                     & Any      & User-specified \\
\bottomrule
\end{tabular}
\begin{tablenotes}\small
\item Note: $\eta = \alpha + \beta t$, $\mu = g^{-1}(\eta)$. The package includes 13 families total (Weibull, Gumbel, Tobit, NegBin, Probit, Beta, ZIP). Custom losses can be supplied via the \texttt{loss=} argument.
\end{tablenotes}
\end{table}

\subsection{Target Functionals}

We are usually not interested in $\theta^*(X_i)$ for every individual, but rather in some summary: an \emph{average} of a function of $\theta^*$ across the population:
\begin{equation}\label{eq:target}
    \mu^* = E[H(X, \theta^*(X), \tilde{t})]
\end{equation}
where $H$ is the \emph{target function} and $\tilde{t}$ is an evaluation point. Common targets include:

\begin{table}[ht]
\centering
\caption{Target functionals for inference.}\label{tab:targets}
\begin{tabular}{lll}
\toprule
\textbf{Target} & $H(x, \theta, \tilde{t})$ & \textbf{Interpretation} \\
\midrule
Average parameter  & $\theta_k$                            & $E[\beta_{\text{price}}]$: mean price sensitivity \\
Average marginal effect & $g'(\alpha + \beta \tilde{t}) \cdot \beta$ & Mean marginal effect at $\tilde{t}$ \\
Elasticity         & $\beta \cdot \tilde{t} \cdot (1-p)$    & Price elasticity of demand \\
Dose-response      & $g(\alpha + \beta \tilde{t})$          & Predicted outcome at $\tilde{t}$ \\
Profit             & $\tilde{t} \cdot g(\alpha + \beta \tilde{t})$ & Expected revenue \\
Custom             & Any $H(x, \theta, \tilde{t})$          & User-specified \\
\bottomrule
\end{tabular}
\end{table}

The Jacobian $H_\theta = \partial H / \partial \theta$ is required for the IF correction. For built-in targets, closed-form Jacobians are provided. For custom targets, the package computes them via automatic differentiation.

\subsection{Why Naive Inference Fails}\label{sec:naive_fails}

To estimate $\mu^*$, we train a neural network $\hat{\theta}(X)$ by minimizing $\sum_i \ell(Y_i, T_i, \hat{\theta}(X_i))$ with sample splitting and then compute $\hat{\mu} = n^{-1}\sum_i H(X_i, \hat{\theta}(X_i), \tilde{t})$.

The \emph{naive} standard error treats $\hat{\theta}$ as if it were the truth and computes $\text{SE}_{\text{naive}} = \text{sd}(H_i) / \sqrt{n}$. This is wrong because $\hat{\theta}$ is itself estimated with error, and neural network regularization introduces systematic bias.

The problem is quantitatively severe. In our simulations (Section~\ref{sec:application} and Appendix~\ref{app:implementation}):
\begin{itemize}
    \item Naive 95\% CIs achieve only 0--20\% actual coverage
    \item The naive SE is 3--10$\times$ too small
    \item The bias is toward zero (regularization shrinks $\hat{\theta}$)
\end{itemize}

\noindent This is not a small-sample problem; it persists at $n = 50{,}000$. The regularization bias is a \emph{feature} of neural network training (it prevents overfitting) but becomes a \emph{bug} for inference.

To see why formally, consider the decomposition of the naive estimator's error:
\begin{equation}\label{eq:bias_decomp}
    \sqrt{n}\!\left(\hat{\mu}_{\text{naive}} - \mu^*\right) = \underbrace{\frac{1}{\sqrt{n}}\sum_{i=1}^n \left[H_i - \mu^*\right]}_{\text{sampling noise}} + \underbrace{\frac{1}{\sqrt{n}}\sum_{i=1}^n H_{\theta,i}\!\left(\hat{\theta}_i - \theta^*_i\right)}_{\text{regularization bias}} + o_P(1)
\end{equation}
The first term is standard sampling variability and vanishes at rate $\sqrt{n}$. The second term captures the bias introduced by the neural network's regularization: because $\hat{\theta}$ is systematically shrunk toward zero (by weight decay, dropout, and early stopping), $\hat{\theta}_i - \theta^*_i$ is not mean-zero and this term does not vanish.\footnote{In our validation suite, naive 95\% CIs achieve 0--20\% coverage; IF-corrected CIs achieve 96\% ($M=50$, $n=5{,}000$, binary logit, eval\_06).} The IF correction in the next section is designed precisely to cancel this second term.

\subsection{The Influence Function Correction}\label{sec:if_correction}

The FLM influence function correction removes the regularization bias. For each observation $i$, define:
\begin{equation}\label{eq:psi}
    \psi_i = H_i - H_{\theta,i} \cdot \Lambda_i^{-1} \cdot \ell_{\theta,i}
\end{equation}
where:
\begin{itemize}
    \item $H_i = H(X_i, \hat{\theta}(X_i), \tilde{t})$ is the target at estimated parameters
    \item $H_{\theta,i} = \frac{\partial H}{\partial \theta}\big|_{\hat{\theta}(X_i)} \in \mathbb{R}^{1 \times d_\theta}$ is the target Jacobian
    \item $\Lambda_i = E[\nabla_\theta^2 \ell(Y, T, \theta) \mid X = X_i]\big|_{\hat{\theta}(X_i)} \in \mathbb{R}^{d_\theta \times d_\theta}$ is the expected Hessian
    \item $\ell_{\theta,i} = \nabla_\theta \ell(Y_i, T_i, \hat{\theta}(X_i)) \in \mathbb{R}^{d_\theta}$ is the score
\end{itemize}

The debiased estimator and its standard error are:
\begin{equation}\label{eq:mu_hat}
    \hat{\mu} = \frac{1}{n}\sum_{i=1}^n \psi_i, \qquad \widehat{\text{SE}} = \frac{\text{sd}(\psi_1, \ldots, \psi_n)}{\sqrt{n}}
\end{equation}

\paragraph{Intuition.} The correction term $H_{\theta} \Lambda^{-1} \ell_\theta$ is an ``adjustment for what the neural net got wrong.'' The score $\ell_\theta$ measures how far observation $i$'s data is from the fitted model. The Hessian $\Lambda$ converts this into parameter-space error. The Jacobian $H_\theta$ maps this to the target. When the neural net fits perfectly (scores near zero), the correction vanishes. When it makes systematic errors (regularization bias), the correction kicks in.

\paragraph{Neyman orthogonality.} The IF $\psi$ in Equation~\eqref{eq:psi} satisfies the \emph{Neyman orthogonality} condition:
\begin{equation}\label{eq:neyman_orth}
    \frac{\partial}{\partial \theta} E[\psi(Y, T, \theta, \Lambda)] \bigg|_{\theta = \theta^*, \Lambda = \Lambda^*} = 0
\end{equation}
The correction term $H_\theta \Lambda^{-1} \ell_\theta$ is constructed precisely so that first-order perturbations in $\hat{\theta}$ around $\theta^*$ cancel. This is why the debiased estimator $\hat{\mu} = \bar{\psi}$ achieves $\sqrt{n}$-consistency despite the neural network's slow convergence rate: the estimation error in $\hat{\theta}$ enters only through a second-order remainder, which is negligible at the $\sqrt{n}$ scale.\footnote{Package autodiff matches closed-form oracle scores and Hessians to machine precision ($< 10^{-15}$) on the multinomial logit (eval\_02), ensuring that the Neyman orthogonality condition is satisfied numerically.}

\paragraph{Connection to classical statistics.} Equation~\eqref{eq:psi} is a semiparametric influence function in the tradition of \citet{hampel1974influence} and \citet{newey1994large}. The specific form arises from the ``one-step correction'' or ``debiasing'' approach used widely in semiparametric econometrics \citep{chernozhukov2018double}. The contribution of FLM is showing that this correction remains valid when the first-stage nuisance (the structural parameters) is estimated by a deep neural network.

\subsection{Three Regimes for Lambda}\label{sec:regimes}

The expected Hessian $\Lambda(X) = E[\nabla_\theta^2 \ell(Y, T, \theta) \mid X]$ depends on how the model and data are structured. Three cases arise:

\begin{table}[ht]
\centering
\caption{Three Lambda estimation regimes.}\label{tab:regimes}
\begin{tabular}{llll}
\toprule
\textbf{Regime} & \textbf{When} & \textbf{Lambda method} & \textbf{Cross-fitting} \\
\midrule
A & RCT, known $F_T$ & Compute (MC integration) & 2-way \\
B & Linear model      & Analytic (closed-form)   & 2-way \\
C & Obs.\ + nonlinear & Estimate (ridge)         & 3-way \\
\bottomrule
\end{tabular}
\end{table}

\paragraph{Regime A: Randomized experiment with known treatment distribution.}
If $T$ is randomly assigned with known distribution $F_T$, and the Hessian does not depend on $Y$, then $\Lambda(X)$ can be computed via Monte Carlo integration:
$\Lambda(X) = \int \nabla_\theta^2 \ell(y, t, \theta) \, dF_T(t)$.
This is the simplest case. Examples: A/B tests, randomized pricing experiments.

\paragraph{Regime B: Linear model.}
For linear models, the Hessian is constant ($\nabla_\theta^2 \ell = 2$), so $\Lambda(X) = 2E[T T' \mid X]$, which can be estimated analytically. This avoids the need for three-way splitting.

\paragraph{Regime C: Observational data with nonlinear model.}
In most applied settings, including our H\&M application, the Hessian depends on $\theta$ (e.g., through $p(1-p)$ in logit or softmax probabilities in multinomial logit). Since $\theta$ is estimated, we cannot simply compute $\Lambda$; we must \emph{estimate} it. FLM propose regressing the individual Hessians on $X$ using a ridge regression. This requires a \emph{three-way} sample split: one part for estimating $\hat{\theta}$, one part for estimating $\hat{\Lambda}$, and one part for assembling the IF. The package handles this automatically.

\subsection{The Multinomial Choice Model}\label{sec:mnl}

Our application uses a multinomial logit (conditional logit) model \citep{mcfadden1974conditional, train2009discrete}. Consumer $i$ chooses from $J$ alternatives, each described by $K$ attributes $x_{ij} \in \mathbb{R}^K$. Utilities are:
\begin{equation}
    V_{ij} = x_{ij}' \theta(X_i), \qquad P(Y_i = j \mid X_i) = \frac{e^{V_{ij}}}{\sum_{m=1}^J e^{V_{im}}}
\end{equation}

This is closely related to the BLP random coefficients model \citep{berry1995automobile, nevo2000practitioner}. The key differences are:
\begin{enumerate}
    \item BLP estimates a distribution (e.g., $\theta \sim N(\bar{\theta}, \Sigma)$) using aggregate market-level data. FLM estimates $\theta(X_i)$ for each individual using individual-level data.
    \item BLP uses GMM with instruments. FLM uses neural networks with influence function correction.
    \item BLP requires solving a fixed-point contraction mapping for each market. FLM requires only a forward pass through the neural network.
\end{enumerate}

The connection between two-tower embeddings and the FLM framework is natural: the consumer embedding $X_i$ captures the latent heterogeneity dimensions, and the neural network maps $X_i \to \theta(X_i)$. The item attributes $T_i$ include prices and (PCA-reduced) item embeddings.

\subsection{Available Models and Targets}

The \texttt{deep-inference} package provides 13 built-in model families (Table~\ref{tab:models}) and 10+ target functionals (Table~\ref{tab:targets}). For models not in the built-in library, users can supply a custom loss function:

\begin{lstlisting}
import torch
from deep_inference import inference

def my_loss(y, t, theta):
    """Custom structural loss."""
    mu = torch.exp(theta[0] + theta[1] * t)  # log-link
    return mu - y * torch.log(mu)             # Poisson NLL

result = inference(Y, T, X, loss=my_loss, theta_dim=2,
                   target_fn=lambda x, th, tt: th[1])
\end{lstlisting}

The package automatically determines whether the Hessian depends on $\theta$ (triggering Regime C and three-way splitting) and computes all required derivatives via automatic differentiation. The user needs only the loss function and the target.
